\section{Phụ lục}
\label{sec:appendix}

\subsection{Công thức độ mạnh xu hướng và mùa}

Với phân rã STL cho chuỗi univariate (hoặc kênh mục tiêu chính), gọi \(T\), \(S\), \(R\)
lần lượt là xu hướng, mùa và phần dư. Theo Hyndman / FPP3 \cite{hyndman2021fpp3}:

\begin{align*}
  F_T &= \max\!\bigl(0,\, 1 - \mathrm{Var}(R)/\mathrm{Var}(T+R)\bigr), \\
  F_S &= \max\!\bigl(0,\, 1 - \mathrm{Var}(R)/\mathrm{Var}(S+R)\bigr).
\end{align*}

Chu kỳ STL chọn theo sampling (ví dụ ETTh theo giờ: \(24\); chuỗi ngày: \(7\) hoặc \(365\)
nếu đủ dài) và sẽ được ghi trong file profile WEEK02.

\subsection{Lưới lookback và horizon dự kiến}

\begin{center}
\small
\begin{tabular}{@{}lll@{}}
\toprule
\textbf{Nhóm} & \textbf{Lookback \(L\)} & \textbf{Horizon \(T\)} \\
\midrule
LTSF chuẩn
  & \(\{96, 192, 336, 720\}\)
  & \(\{96, 192, 336, 720\}\) \\
\addlinespace
VIC (Vingroup ngày)
  & \(\{5, 30, 120, 480\}\)
  & \(5\) \\
\bottomrule
\end{tabular}
\end{center}

Nếu thiếu thời gian tính toán, lưới LTSF có thể cắt còn \(L \in \{96,336\}\) và
\(T \in \{96,336\}\); mọi cắt giảm phải ghi trong protocol WEEK03.

\subsection{Đường dẫn Working Files liên quan}

\begin{center}
\small
\begin{tabular}{@{}>{\raggedright\arraybackslash}p{0.32\textwidth}
                 >{\raggedright\arraybackslash}p{0.58\textwidth}@{}}
\toprule
\textbf{Nội dung} & \textbf{Đường dẫn} \\
\midrule
Outline (RQ, quyết định)
  & \code{Working_Files/WEEK01_Outline.md} \\
\addlinespace
Reading note (cơ chế)
  & \code{Working_Files/WEEK01_Reading_Note/} \\
\addlinespace
Bài research-gap này
  & \code{Working_Files/WEEK02 - Research Gap Analysis/} \\
\addlinespace
Paper Tracker
  & \code{Paper_Tracker/} \\
\addlinespace
Kế hoạch thực thi
  & \code{EXECUTION_PLAN.md} \\
\bottomrule
\end{tabular}
\end{center}

\subsection{Những gì WEEK02 chưa có}

Tài liệu này chưa chứa bảng MSE/MAE thực nghiệm, scatter RQ1, hay đường cong RQ2.
Các phần đó thuộc WEEK03 (baseline) và WEEK04 (ablation + kết luận), và sẽ được gắn thêm
vào cùng file LaTeX này.
